\chapter*{RESUMEN}
\addcontentsline{toc}{chapter}{RESUMEN}

La gestión del soporte técnico y el control de activos informáticos para empresas clientes externas en J\&P Periféricos S.A.C. se veía afectada por la recepción de solicitudes mediante llamadas y WhatsApp sin registro formal, el uso de archivos de Excel dispersos para los expedientes de equipos, cruces de agenda en visitas \textit{in-situ}, revisiones manuales de contratos y SLA, y la falta de indicadores de desempeño. Este escenario generaba demoras operativas, reprocesos, pérdidas económicas estimadas y una menor trazabilidad en la atención brindada a cada cliente, sucursal y trabajador.

Ante ello, el proyecto tuvo como objetivo desarrollar e implementar un Sistema Web Help Desk Centralizado orientado a las empresas clientes externas de J\&P Periféricos S.A.C., con el fin de ordenar la atención de incidencias, asegurar la trazabilidad del inventario de hardware y software y elevar la satisfacción del servicio. Para el diagnóstico de las causas raíz (CA1 a CA7), se aplicó un enfoque cuantitativo con diseño preexperimental a partir de entrevistas, revisión de registros operativos y análisis de Pareto sobre 15 evaluadores. En el desarrollo se integraron prácticas de ITSM, ITIL v4 e ITAM sobre una arquitectura modular compuesta por un \textit{backend} en NestJS con Prisma ORM sobre MySQL, MongoDB para el historial de chats, Redis (\textit{Pub/Sub}) y WebSockets para eventos en tiempo real, y un \textit{frontend} en Vite, React y \textit{shadcn/ui} estructurado por roles.

La plataforma implementada centraliza los tickets con código PIN único y flujo de estados auditable, organiza el expediente e inventario multinivel por empresa, sucursal, área y trabajador, coordina visitas \textit{in-situ} evitando cruces horarios, valida planes y SLA de forma digital, estandariza el escalamiento y consolida tableros de KPIs (MTTA, MTTR, FCR y CSAT). Con ello, se reemplazan los canales informales, se reducen las búsquedas manuales y se respalda la toma de decisiones del área. Se concluye que el sistema atiende de manera integral la problemática diagnosticada y fortalece el control y la calidad del servicio de soporte técnico de J\&P Periféricos S.A.C.

Palabras clave: Sistema web Help Desk, gestión de soporte técnico, control de inventario de activos TI, arquitectura hexagonal, trazabilidad de incidencias.